\section{Local cyclicity of \texorpdfstring{\(I_2^1\)}{I21}}
\label{sec:introduction}

The Dumortier--Roussarie--Rousseau program reduces the finiteness
problem for quadratic limit cycles to local cyclicity questions for
specified limit-periodic sets~\cite{DRR1994}. We study the loop
\(I_2^1\) through a nilpotent saddle--node at infinity. Our aim is
to count the cycles approaching the entire loop, including parameter
directions in which boundary roots merge or coefficient scales vanish
independently, and to attain the resulting bound in the original family.

Fix \(C_0>1\). We consider the five-parameter quadratic family
\begin{equation}
\label{eq:quadratic-family}
\begin{aligned}
 \dot x&=Ax-y+x^2+(\mu_2+\mu_3)xy+\mu_1y^2,\\
 \dot y&=Cx+x^2+xy+\mu_3y^2
\end{aligned}
\end{equation}
near the parameter
\begin{equation}
\label{eq:organizing-parameter}
 (A,C,\mu_1,\mu_2,\mu_3)=(1,C_0,0,0,0).
\end{equation}
At this parameter the invariant parabola
\begin{equation}
\label{eq:organizing-parabola}
 y=(x^2-C_0)/2
\end{equation}
and its resolved nilpotent saddle--node passage at infinity form
the directed graphic \(\GammaI\), denoted \(I_2^1\) in the DRR
classification \cite[Proposition~4.2.2(6), Figure~8(6)]{DRR1994}
(see \cref{sec:graphic-central-return,fig:organizing-graphic-new}).

We measure proximity to this graphic on the Poincar\'e sphere
\(\mathbb S^2\), where the Poincar\'e compactification of
\eqref{eq:quadratic-family} is an analytic vector field. Let
\(\distH\) denote Hausdorff distance for a fixed metric on
\(\mathbb S^2\). When resolved or scaled coordinates are used below,
the sets are projected back to \(\mathbb S^2\) before their Hausdorff
distance is taken.

For \(\varepsilon>0\) and a neighborhood \(V\subset\mathbb R^5\) of the
organizing parameter \eqref{eq:organizing-parameter}, let
\(N(V,\varepsilon)\) be the supremum over \(p\in V\) of the number of
limit cycles \(\gamma\) of \eqref{eq:quadratic-family} with parameter
\(p\) such that \(\distH(\overline\gamma,\GammaI)<\varepsilon\).
Following the Hausdorff formulation of cyclicity in
\cite[Definition~2.2(2)]{DRR1994} (see also
\cite[\S2.2.1]{Roussarie1998}), the cyclicity of
\(\GammaI\) in the family \eqref{eq:quadratic-family} is
\begin{equation}
\label{eq:cyclicity-definition}
 \Cycl(\GammaI)=\inf_{V,\,\varepsilon}N(V,\varepsilon).
\end{equation}
The Hausdorff condition requires proximity to the whole graphic.

\begin{theorem}[Local upper bound]
\label{thm:direct-finite-cyclicity}
For every fixed \(C_0>1\), there exist a neighborhood
\(V\subset\mathbb R^5\) of \((1,C_0,0,0,0)\) and a number
\(\varepsilon>0\) such that, for every \(p\in V\), the vector field
\eqref{eq:quadratic-family} with parameter \(p\) has at most three
periodic orbits \(\gamma\) with
\[
 \distH(\overline\gamma,\Gamma(C_0))<\varepsilon .
\]
Each of these periodic orbits is a limit cycle. In particular,
\(\Cycl(\GammaI)\le3\).
\end{theorem}

\begin{proposition}[Three periodic orbits]
\label{prop:three-cycle-lower-bound}
For every fixed \(C_0>1\), there is a sequence of parameters
\(p_n\to(1,C_0,0,0,0)\), with \(C=C_0\), for which
\eqref{eq:quadratic-family} has three distinct isolated periodic
orbits \(\gamma_{j,n}\), \(j=1,2,3\), satisfying
\[
 \distH(\overline{\gamma_{j,n}},\Gamma(C_0))\longrightarrow0.
\]
These orbits eventually lie in every fixed sufficiently small
neighborhood of the whole graphic, and in particular satisfy the
Hausdorff condition of \cref{thm:direct-finite-cyclicity}.
\end{proposition}

\begin{corollary}[Cyclicity three]
\label{cor:exact-local-cyclicity-three}
For every fixed \(C_0>1\), \(\Cycl(\GammaI)=3\).
\end{corollary}

The three orbits of \cref{prop:three-cycle-lower-bound} are also
hyperbolic and persist on open sets in the original five-parameter space
arbitrarily close to the organizing parameter
(\cref{cor:hyperbolic-triple}).

The neighborhood \(V\) and the radius \(\varepsilon\) depend on
\(C_0\). Their construction uses, in particular, the positive
speed \eqref{eq:parabola-speed-new} of the organizing orbit and the
limiting contraction multiplier \(\Lambda_0<1\) of
\cref{lem:exterior-multiplier-new}.
The speed bound \((C_0-1)/2\) tends to zero as \(C_0\downarrow1\),
and the margin \(1-\Lambda_0\) tends to zero as \(C_0\to\infty\).

\subsection*{Relation to previous work}

The DRR program develops normal forms near saddle--nodes and applies them to
finite-cyclicity questions; see Dumortier, Ilyashenko, and Rousseau
\cite{DIR2002}. Related parts of the program treat the degenerate graphics
$DF_{1a}$ and $DF_{2a}$ \cite{DR2009}, with the latter study completed in the
remaining chart by Huzak \cite{Huzak2018}. For graphics through
triple nilpotent points, Roussarie and Rousseau study the center graphic
$I_{14}^1$ and distinguish it from the blown-up boundary limit-periodic sets
appearing in the same analysis \cite{RR2015}; Rousseau, Shan, and Zhu treat
the saddle-type graphics $I_{12}^1$ and $I_{13}^1$ \cite{RSZ2016}.

Huzak and Kristiansen develop entry--exit formulas for degenerate turning
points and apply them to the present five-parameter family
\cite[Theorem~2.4 and \S6]{HK2026}. Their theorem constructs a local Dulac
map, with finite-order remainder estimates in a logarithmic parameterization,
in a principal root-free region where the scaled boundary drift has a uniform
sign margin. Their Section~6 derives the entry--exit relation for this
family and leaves the cyclicity analysis to further work. The passage
problem here also includes merging boundary
roots, roots approaching the central section, and independently
vanishing coefficient scales. We compare the two endpoint maps on the
full interval of admissible central heights, compose them with the
exterior transition, and use geometric capture to count the periodic
orbits satisfying the whole-graphic Hausdorff condition.

The root-free calculation enters our buffered compact-exit analysis
through the leading relation between normalized exit radii,
\[
 f_\Lambda(\rho)=\frac{\Lambda\rho}{1+(1+\Lambda)\rho}.
\]
Here \(\rho\) is the normalized positive exit radius and \(\Lambda\)
is the boundary-action multiplier. This relation is the inverse,
after a sign change and positive rescaling, of the base-point relation
in Huzak--Kristiansen \cite[(2.13) and (6.8)]{HK2026}.
The change of variables is given after
\eqref{eq:direct-buffered-fractional-relation}.

Uniform \(C^3\) matching transfers this relation to the endpoint maps
used for counting. Drift-weighted estimates treat the other passage
regions, and a scattering comparison treats contact with the central
section. The parameter cover combines these comparisons into the
upper bound. A direct matching construction and variation of the
connection splitting then attain it with \(C=C_0\).

Finite-order passage analysis for canard cycles with singular slow dynamics,
including divergence comparisons and exponential-flatness estimates, was
developed by De Maesschalck and Dumortier \cite{DMD2008}. For limit-periodic
sets associated with singular slow dynamics, their later work obtains
cycle bounds by studying suitable derivatives of the slow-divergence integral,
even when that integral is unbounded \cite{DMD2011}.
Mar\'in and Villadelprat give parameter-dependent
hyperbolic-saddle Dulac expansions, including the near-resonant case
\cite{MV2024}. These works provide the background for the passage
estimates and derivative comparisons developed here.

\subsection*{Proof strategy}

The proof has two counting steps. For fixed parameters, we control
one return displacement \(\Delta_p\) on its full interval
\(\mathcal H_p\) of central heights. A finite cover of parameter
space then selects an applicable estimate throughout a neighborhood
of the organizing value. Geometric capture connects these zeros to
the periodic orbits in the theorem.

\emph{A common height for the two halves.}
In normal-rate coordinates \((\eta,z)\), the field is
\[
 \dot\eta=q(\eta,p)-zB(\eta,z,p),\qquad \dot z=-\eta z.
\]
Here \(q(\eta,p)=q_0+q_1\eta+q_2\eta^2+\eta^3R_q(\eta,p)\)
is the boundary drift. At fixed nearby \(A,C\), its coefficients
\((q_0,q_1,q_2)\) are local coordinates in
\((\mu_1,\mu_2,\mu_3)\) (Appendix~\ref{app:coordinates}).
The central height \(h\) is the coordinate \(z\) on \(\eta=0\).
From \((0,h)\), follow the positive half-orbit backward and the negative
half-orbit forward. Their endpoint functions \(I_\pm(h)\) are defined
on the same interval \(\mathcal H_p\). The regular exterior map
\(\Pi_p\) maps the negative endpoint window to the positive one.
The return displacement is
\[
 \Delta_p(h)=\Pi_p(I_-(h))-I_+(h).
\]
A zero expresses the matching of the two endpoints through the exterior
transition.

\emph{From endpoint derivatives to a zero count.}
Both endpoint derivatives with respect to \(h\) are positive. The
logarithmic slope ratio
\(\mathcal L=\log[\Pi_p'(I_-)I_{-,h}/I_{+,h}]\) satisfies
\[
 \Delta_{p,h}=I_{+,h}(\e^{\mathcal L}-1).
\]
The zeros of \(\mathcal L\) are therefore the critical points of
\(\Delta_p\). To bound them on the full height interval, we combine
crossing directions on the outer height ranges with negative
curvature at critical points on the middle range. The order of these
ranges allows at most two zeros of \(\mathcal L\); Rolle's theorem
then gives at most three zeros of \(\Delta_p\).

\emph{Uniformity near degenerate passages.}
The comparison must remain uniform when roots degenerate and coefficient
scales vanish independently. Each endpoint derivative retains the tail
multiplier \((z_*/t)^{-2q_2}\), even though the common quadratic term
cancels from their ratio; this allows unbounded \(q_2\log t\).
Near a boundary root, the leading curvature vanishes with the drift.
The boundary-flow coordinate of \cref{lem:direct-boundary-flow}
absorbs the height-dependent radial term, and
\cref{lem:direct-weighted-phase} makes the phase correction smaller
than any prescribed finite drift power after a fixed positive
boundary-time advance. The resulting drift-weighted curvature formula
\eqref{eq:direct-rate-curvature} preserves every fixed sign margin.

The exterior transition supplies the remaining contraction when the
singular slopes nearly agree. Its organizing derivative tends to
\[
 \Lambda_0=\exp\!\left(-\frac{4\pi}{\sqrt{C_0-1}}\right)<1
\]
as the sections approach infinity. Fixed sections retain a contraction
margin, while the nonconstant exterior term carries the same drift
weights as the passage terms. At central contact, a decreasing scattering
derivative ratio with a small \(C^2\) error gives at most one middle
critical point. The outer derivative signs then give at most two
contact zeros on the full height interval.

\emph{Realizing three cycles.}
For the lower bound, keep \(C=C_0\) and choose a coefficient path in
the root-free region with negative quadratic drift. Comparing the two
boundary actions gives a compact exit profile with a strict negative
minimum. The parameter \(A\) leaves the boundary drift unchanged and
splits the organizing connection transversely, so it can adjust the
additive constant of this profile. Three test exits then give two
displacement zeros.

This compact exit comparison alone does not locate the third cycle.
Returning to the common central height, the organizing contraction gives
the opposite displacement sign at a larger height. This supplies a third
zero beyond the heights tested by the compact profile. Graph order and
regular continuation then show that all three orbits converge to the
whole graphic. The coefficient path, action matching, and splitting
calculation are developed in \cref{sec:sharpness}; hyperbolicity and the
separation of the third height follow in
\cref{cor:hyperbolic-triple,cor:third-height-separation}.

The proof is organized around these comparisons. The geometry and the
common-height return problem are established in
\cref{sec:graphic-central-return}, and
\cref{sec:standard-tools,sec:direct-passage-estimates} give the counting
principles and uniform passage estimates. Contact and positive offsets
are treated in \cref{sec:contact-positive-offsets}. For negative offsets,
\cref{sec:recurring-ranges} develops the height-range comparisons used on
both coefficient scales; the primary-scale analysis is in
\cref{sec:primary-regimes}, and the secondary-scale cases begin in
\cref{sec:quadratic-regimes}. These are estimates on overlapping height
intervals for fixed parameters. Their coverage of a full parameter
neighborhood is proved separately in \cref{sec:direct-parameter-cover},
where \cref{tab:direct-cover-guide} identifies the applicable comparison
and bound in each region.
